% Part: second-order-logic
% Chapter: sol-and-set-theory
% Section: power-of-continuum

\documentclass[../../../include/open-logic-section]{subfiles}

\begin{document}

\olfileid{sol}{set}{pow}

\olsection{توان پیوستار}

\begin{explain}
در منطق مرتبهٔ دوم می‌توانیم بر زیرمجموعه‌های !!{domain} کمیت‌گذاری
کنیم، اما نمی‌توانیم بر مجموعه‌های زیرمجموعه‌های !!{domain} کمیت‌گذاری
کنیم. برای انجام مستقیم این کار به منطق \emph{مرتبهٔ سوم} نیاز داشتیم.
برای نمونه، اگر می‌خواستیم قضیهٔ کانتور را بیان کنیم که هیچ تابع
!!{injective}ی از مجموعهٔ توانیِ یک مجموعه به خود آن مجموعه وجود ندارد،
شاید می‌کوشیدیم آن را چنین صورت‌بندی کنیم: «برای هر مجموعهٔ~$X$ و هر
مجموعهٔ~$P$، اگر~$P$ مجموعهٔ توانی~$X$ باشد، آنگاه نه
$\cardle{P}{X}$». برای گفتن اینکه~$P$ مجموعهٔ توانی~$X$ است، باید این
مطلب را صورت‌بندی کنیم که !!{element}s مجموعهٔ~$P$ دقیقاً همهٔ
زیرمجموعه‌های~$X$ هستند؛ یعنی چیزی مانند
$\lforall[Y][(P(Y) \liff Y \subseteq X)]$. مشکل در~$P(Y)$ نهفته است:
این عبارت !!a{formula}ی از منطق مرتبهٔ دوم نیست، زیرا تنها ترم‌ها می‌توانند
آرگومان متغیرهای رابطه‌ای یک‌موضعی چون~$P$ باشند.

اما اگر دامنه به‌اندازهٔ کافی بزرگ باشد، می‌توانیم کمیت‌گذاری بر
مجموعه‌های مجموعه‌ها را \emph{شبیه‌سازی} کنیم. ایده این است که از این
واقعیت بهره بگیریم که رابطهٔ دو‌موضعی~$R$، !!{element}s
!!{domain} را با !!{element}s !!{domain} مرتبط می‌کند. با داشتن چنین
$R$ای می‌توانیم همهٔ !!{element}sی را گرد آوریم که~$x$ با آن‌ها رابطهٔ
$R$ دارد: $\Setabs{y \in \Domain{M}}{R(x,y)}$ مجموعه‌ای است که~$x$ آن
را «رمزگذاری می‌کند». برعکس، اگر~$Z \subseteq \Pow{\Domain{M}}$
مجموعه‌ای از زیرمجموعه‌های~$\Domain{M}$ باشد و !!{element}s
$\Domain{M}$ دست‌کم به‌اندازهٔ مجموعه‌های~$Z$ باشند، آنگاه رابطه‌ای چون
$R \subseteq \Domain{M}^2$ نیز وجود دارد که هر~$Y \in Z$ را عضوی
چون~$x$ با استفاده از~$R$ رمزگذاری می‌کند.
\end{explain}

\begin{defn}
اگر~$R \subseteq \Domain{M}^2$، آنگاه \emph{$x$ به‌کمک~$R$}
$\Setabs{y \in \Domain{M}}{R(x, y)}$ را \emph{رمزگذاری می‌کند}.
\end{defn}

اگر !!a{element}ی چون~$x \in \Domain{M}$ به‌کمک~$R$ مجموعه‌ای چون
  $Z \subseteq \Domain{M}$ را رمزگذاری کند، آنگاه مجموعه‌ای
چون~$Y \subseteq \Domain{M}$ مجموعه‌ای از مجموعه‌ها، یعنی مجموعه‌های
رمزگذاری‌شده به‌وسیلهٔ !!{element}s مجموعهٔ~$Y$، را رمزگذاری می‌کند.
پس~$Y$ می‌تواند به‌کمک~$R$، $\Pow{X}$ را رمزگذاری کند. این امر اگر و
تنها اگر برای هر~$Z \subseteq X$ عضوی چون~$x \in Y$ به‌کمک~$R$، $Z$
را رمزگذاری کند و هر~$x \in Y$ به‌کمک~$R$، مجموعه‌ای چون
$Z \subseteq X$ را رمزگذاری کند.

\begin{prop}
!!{formula}
  \[
  \fn{Codes}(x, R, Z) \ident \lforall[y][(Z(y) \liff
    R(x, y))]
  \]
بیان می‌کند که~$s(x)$ به‌کمک~$s(R)$، $s(Z)$ را رمزگذاری می‌کند.
!!{formula}
\begin{multline*}
  \fn{Pow}(Y, R, X) \ident \\
  \lforall[Z][(Z \subseteq X \lif \lexists[x][(Y(x) \land
    \fn{Codes}(x, R, Z))])] \land {} \\ \lforall[x][(Y(x) \lif
  \lforall[Z][(\fn{Codes}(x, R, Z) \lif Z \subseteq X)])]
\end{multline*}
بیان می‌کند که~$s(Y)$ به‌کمک~$s(R)$، مجموعهٔ توانی~$s(X)$ را رمزگذاری
می‌کند؛ یعنی اعضای~$s(Y)$ به‌کمک~$s(R)$ دقیقاً زیرمجموعه‌های~$s(X)$ را
رمزگذاری می‌کنند.
\end{prop}

\begin{explain}
با این شگرد می‌توانیم با کمیت‌گذاری بر رمزهای زیرمجموعه‌ها، به‌جای خود
زیرمجموعه‌ها، گزاره‌هایی دربارهٔ مجموعهٔ توانی بیان کنیم. برای نمونه،
اکنون می‌توان قضیهٔ کانتور را با این گفته بیان کرد که هیچ تابع
!!{injective}ی از دامنهٔ هیچ رابطه‌ای که مجموعهٔ توانی~$X$ را رمزگذاری
می‌کند به خود~$X$ وجود ندارد.
\end{explain}

\begin{prop}
جملهٔ
\begin{multline*}
  \lforall[X][\lforall[Y][\lforall[R][(\fn{Pow}(Y, R, X) \lif \\
      \lnot \lexists[u][(\lforall[x][\lforall[y][(\eq[u(x)][u(y)] \lif
            \eq[x][y])]] \land {}\\
        \lforall[x][(Y(x) \lif X(u(x)))])])]]]
\end{multline*}
معتبر است.
\end{prop}

\begin{explain}
مجموعهٔ توانیِ مجموعه‌ای~!!a{denumerable}، !!{nonenumerable} است، پس توان
آن از توان هر مجموعهٔ !!{denumerable}ی (یعنی~$\aleph_0$) بزرگ‌تر است.
اندازهٔ~$\Pow{\Nat}$ «توان پیوستار» نامیده می‌شود، زیرا با اندازهٔ نقاط
خط اعداد حقیقی، یعنی~$\Real$، برابر است. اگر !!{domain} به‌اندازهٔ کافی
بزرگ باشد تا مجموعهٔ توانیِ مجموعه‌ای~!!a{denumerable} را رمزگذاری کند،
می‌توانیم با گفتن اینکه یک مجموعه با هر مجموعهٔ~$Y$ای که مجموعهٔ توانیِ
مجموعهٔ~$X$ با اندازهٔ~$\aleph_0$ را رمزگذاری می‌کند هم‌شمار است،
اندازهٔ پیوستارداشتن آن را بیان کنیم. (اگر دامنه به‌اندازهٔ کافی بزرگ
نباشد، یعنی هیچ زیرمجموعهٔ هم‌شمار با~$\Real$ نداشته باشد، آنگاه هیچ
رابطه‌ای نیز نمی‌تواند~$\Pow{X}$ را رمزگذاری کند.)
\end{explain}

\begin{prop}
اگر~$\cardle{\Real}{\Domain{M}}$، آنگاه !!{formula}
\begin{multline*}
\fn{Cont}(Y) \ident
\lexists[X][\lexists[R][((\fn{Aleph}_0(X) \land
      \fn{Pow}(Y, R, X)) \land \\
      \lforall[x][\lforall[y][((Y(x) \land {}
      Y(y) \land \lforall[z][R(x,z) \liff R(y,z)]) \lif \eq[x][y])]])]]
\end{multline*}
بیان می‌کند که~$\cardeq{s(Y)}{\Real}$.
\end{prop}

\begin{proof}
  $\fn{Pow}(Y, R, X)$ بیان می‌کند که~$s(Y)$ به‌کمک~$s(R)$ مجموعهٔ
  توانی~$s(X)$ را رمزگذاری می‌کند و~$\fn{Aleph}_0(X)$ شمارابودن آن را
  بیان می‌کند. پس~$s(Y)$ دست‌کم به‌اندازهٔ توان پیوستار بزرگ است، هرچند
  ممکن است بزرگ‌تر باشد (اگر چند عضو~$s(Y)$ زیرمجموعهٔ یکسانی از~$X$ را
  رمزگذاری کنند). عطف آخر این امکان را منتفی می‌کند، زیرا می‌خواهد
  تناظر میان اعضای~$s(Y)$ و زیرمجموعه‌های~$s(X)$ به‌واسطهٔ~$s(R)$
  !!{injective} باشد.
\end{proof}

\begin{prop}
$\cardeq{\Domain{M}}{\Real}$ اگر و تنها اگر
\begin{multline*}
  \Sat{M}{\lexists[X][\lexists[Y][\lexists[R][(\fn{Aleph}_0(X) \land \fn{Pow}(Y, R, X) \land \\
      \lforall[x][\lforall[y][((Y(x) \land Y(y) \land \lforall[z][R(x,z) \liff R(y,z)]) \lif \eq[x][y])]] \land {} \\
      \lexists[u][(\lforall[x][\lforall[y][(\eq[u(x)][u(y)] \lif \eq[x][y])]] \land {} \\
      \lforall[x][Y(u(x))] \land {} \\ \lforall[y][(Y(y) \lif \lexists[x][\eq[y][u(x)]])])])]]]}.
        \end{multline*}
  \end{prop}

\begin{explain}
فرضیهٔ پیوستار می‌گوید اندازهٔ پیوستار نخستین توان !!{nonenumerable}
است؛ یعنی~$\Pow{\Nat}$ اندازهٔ~$\aleph_1$ دارد.
\end{explain}

\begin{prop}
فرضیهٔ پیوستار صادق است اگر و تنها اگر
\[\fn{CH} \ident
\lforall[X][(\fn{Aleph}_1(X) \liff \fn{Cont}(X))]\]
معتبر باشد.
\end{prop}

توجه کنید که چنین نیست که~$\lnot \fn{CH}$ معتبر باشد اگر و تنها اگر
فرضیهٔ پیوستار کاذب باشد. در دامنه‌ای~!!a{enumerable} هیچ زیرمجموعه‌ای با
اندازهٔ~$\aleph_1$ و نیز هیچ زیرمجموعه‌ای با اندازهٔ پیوستار وجود ندارد؛
پس~$\fn{CH}$ در دامنه‌ای~!!a{enumerable} همواره صادق است. اما می‌توانیم
جملهٔ دیگری به دست دهیم که معتبر است اگر و تنها اگر فرضیهٔ پیوستار
کاذب باشد:

\begin{prop}
فرضیهٔ پیوستار کاذب است اگر و تنها اگر
\[\fn{NCH} \ident
\lforall[X][(\fn{Cont}(X) \lif \lexists[Y][(Y \subseteq X \land \lnot
    \fn{Count}(Y) \land \lnot \cardeq{X}{Y})])]\]
معتبر باشد.
\end{prop}

\end{document}
